\section{From field error to observable path error}\label{sec:main}
We use the convention
\[
 \TV(P,Q)=\sup_{B}|P(B)-Q(B)|\in[0,1].
\]
For finite sample spaces it equals $\frac12\sum_\omega|P(\omega)-Q(\omega)|$.
The factor $\frac12$ is essential for the execution constants below.

\begin{assumption}[A common admissible simulator]\label{ass:simulator}
The true and learned models have the same initial observed state, the same finite mark set, and the same deterministic transition rule $\Delta$. A common predictable policy $a_t=a(t,\Hist_{t-})$ supplies actions. On the finite horizon $[0,T]$, both observable intensity functionals are measurable and bounded by a finite deterministic envelope. Their point-process constructions are unique in law. The learned functional is defined on every true history needed in the comparison, including histories absent from training.
\end{assumption}

A bounded envelope is a convenient sufficient condition for nonexplosion and common thinning. The comparison extends to nonexplosive locally bounded models by stopping and integrable limits, but the bounded version is all that is required here. A randomized policy is covered by adding a shared random seed to the initial observed information.

\begin{lemma}[Cell integration controls event-rate error]\label{lem:cells}
For $f,g\in L^2(A,\nu)$, disjoint measurable cells $(B_e)$, and common masks $m_e\in[0,1]$,
\begin{equation}\label{eq:cell-contraction}
 \sum_e\left|m_e\int_{B_e}(f-g)\dd\nu\right|
 \le\|f-g\|_{L^1(\nu)}
 \le\sqrt M\,\|f-g\|_{L^2(\nu)}.
\end{equation}
The constant does not depend on the number of cells.
\end{lemma}
\begin{proof}
The absolute integral on each cell is bounded by the integral of the absolute value. Summing over disjoint cells and using $m_e\le1$ gives the first inequality. Cauchy--Schwarz applied to $|f-g|$ and the constant function $1$ gives the second.
\end{proof}

\begin{lemma}[Latent uncertainty and conditional fitting]\label{lem:posterior}
Fix time, observed history, and action. Let true and learned rates be obtained from fields $f(z),\widehat f(z)$ by the same cell integration and mask. Suppose
\begin{equation}\label{eq:field-lip}
 \|\widehat f(z)-\widehat f(z')\|_{L^2(\nu)}
 \le K\|z-z'\|_{\mathsf H}.
\end{equation}
Let $\pi,\widehat\pi\in\Ptwo(\mathsf H)$ and define
\[
 r=\int\|f(z)-\widehat f(z)\|_{L^2(\nu)}\,\pi(\dd z),
 \qquad \kappa=W_2(\pi,\widehat\pi).
\]
Then the posterior-averaged observable intensities satisfy
\begin{equation}\label{eq:rate-error}
 \sum_e|\bar\lambda_e-\widehat{\bar\lambda}_e|
 \le\sqrt M\,(r+K\kappa).
\end{equation}
\end{lemma}
\begin{proof}
Insert the intermediate rate obtained by averaging $\widehat f$ under $\pi$.
For the first difference, integrate \eqref{eq:cell-contraction} under $\pi$, obtaining $\sqrt M r$.
For the second, let $(Z,\widehat Z)$ have any coupling of $\pi,\widehat\pi$. Jensen's inequality for absolute value, \Cref{lem:cells}, and \eqref{eq:field-lip} give
\[
 \sum_e\left|\E\widehat\lambda_e(Z)-\E\widehat\lambda_e(\widehat Z)\right|
 \le\sqrt M K\,\E\|Z-\widehat Z\|
 \le\sqrt M K\,(\E\|Z-\widehat Z\|^2)^{1/2}.
\]
Take the infimum over couplings. The result only needs $W_1$; using $W_2$ permits the orthogonal posterior certificate of the next section.
\end{proof}

\begin{theorem}[Path comparison through observable intensities]\label{thm:path}
Under \Cref{ass:simulator}, let
\begin{equation}
 \delta_t(\Hist)=\sum_e
 |\bar\lambda_e(t,\Hist,a_t)-\widehat{\bar\lambda}_e(t,\Hist,a_t)|.
\end{equation}
The observable path laws, including the event times, marks, and displayed states, satisfy
\begin{equation}\label{eq:path-tv}
 \TV(\mathbb P^a_{[0,T]},\widehat{\mathbb P}^a_{[0,T]})
 \le \min\left\{1,\E^a\int_0^T\delta_t(\Hist_{t-})\dd t\right\}.
\end{equation}
If $\delta_t(\Hist)\le b(t)$ for a deterministic nonnegative integrable function on every common history, then also
\begin{equation}\label{eq:uniform-path-tv}
 \TV(\mathbb P^a_{[0,T]},\widehat{\mathbb P}^a_{[0,T]})
 \le 1-\exp\left(-\int_0^T b(t)\dd t\right).
\end{equation}
\end{theorem}
\begin{proof}
For each mark use a common Poisson random measure on time and a nonnegative auxiliary coordinate, accepting a point when its auxiliary coordinate is below the respective intensity. Continue both constructions after they disagree, using their own histories. Each marginal has its prescribed law.

Let $\sigma$ be the first time the accepted marked histories differ. Before $\sigma$ the observed states, actions, and histories coincide. For mark $e$, disagreement occurs when the auxiliary coordinate lies between the two intensities, an interval of length
$|\bar\lambda_e-\widehat{\bar\lambda}_e|$. Consequently, the first-disagreement counting process has compensator
\[
 \int_0^{T\wedge\sigma}\delta_t(\Hist_{t-})\dd t.
\]
It follows that
\[
 \mathbb P(\sigma\le T)
 =\E\int_0^T\mathbf1_{\{t<\sigma\}}\delta_t(\Hist_{t-})\dd t
 \le\E^a\int_0^T\delta_t(\Hist_{t-})\dd t.
\]
After dropping the indicator, the expectation is under the true marginal. The probability that two coupled paths differ bounds their total variation, proving \eqref{eq:path-tv}.

For the deterministic envelope, put $S(t)=\mathbb P(\sigma>t)$. The compensator identity gives an absolutely continuous survival function with
$-S'(t)=\E[\mathbf1_{\{t<\sigma\}}\delta_t]\le b(t)S(t)$ almost everywhere.
The integrating-factor inequality yields
$S(T)\ge\exp(-\int_0^T b)$, proving \eqref{eq:uniform-path-tv}.
\end{proof}

No metric regularity of $\Delta$ enters the proof. A discontinuous price change after depletion is harmless on the event that the marked histories coincide. Conversely, the result does not compare two different matching engines: if the same mark causes different updates, agreement of intensities alone is insufficient.

\begin{theorem}[The generative-operator market certificate]\label{thm:main}
Under \Cref{ass:simulator}, assume the field representation and hypotheses of \Cref{lem:posterior} hold predictably along true histories, with finite
$\E^a\int_0^T\sqrt M(r_t+K_t\kappa_t)\dd t$. Then
\begin{equation}\label{eq:main}
 \TV(\mathbb P^a_{[0,T]},\widehat{\mathbb P}^a_{[0,T]})
 \le\min\left\{1,\sqrt M\,\E^a\int_0^T(r_t+K_t\kappa_t)\dd t\right\}.
\end{equation}
If implemented observable rates have an additional $\ell^1$ error
$q_t(\Hist)$ relative to the rates just described, add
$\E^a\int_0^T q_t(\Hist_{t-})\dd t$ inside the minimum.
\end{theorem}
\begin{proof}
Insert \eqref{eq:rate-error} into \eqref{eq:path-tv}. For the implemented rates use the triangle inequality for their $\ell^1$ difference before applying the same comparison.
\end{proof}

The extra term can include numerical quadrature, posterior Monte Carlo, and other rate-evaluation errors only when an appropriate bound for the actual predictable implementation is available. A discrete-time sampler that rounds all event times to a grid can have a path law singular to a continuous-time event law. Such time rounding is not automatically represented by a small rate perturbation.

\subsection{Exact posterior randomization at thinning proposals}\label{sec:posterior-thinning}
There is a direct implementation of the averaged intensity that does not hold a sampled latent rate during a waiting interval. It uses fresh randomness at every proposal, including proposals that are rejected.

\begin{proposition}[Posterior-randomized thinning]\label{prop:posterior-thinning}
For each mark $e$, let $c_e>0$ be a deterministic bound for
$\widehat\lambda_e(t,h,z,a)$ on every required time, history, action, and latent state. Generate independent Poisson proposals of rate $c_e$. At a proposal time $t$, evaluate the history-dependent kernel $\widehat\pi_{t-}(h)$, draw a fresh $Z\sim\widehat\pi_{t-}(h)$ and independent $V\sim\mathrm{Unif}(0,1)$, and accept $e$ when
\begin{equation}\label{eq:posterior-thinning}
 V\le\widehat\lambda_e(t,h,Z,a)/c_e.
\end{equation}
The accepted process has observable intensity
$\int\widehat\lambda_e(t,h,z,a)\widehat\pi_{t-}(h,\dd z)$.
The same conclusion holds with one proposal stream of rate $c$ whenever
$\sum_e\widehat\lambda_e\le c$, by selecting mark $e$ with probability
$\widehat\lambda_e/c$ and rejecting the remaining probability.
\end{proposition}
\begin{proof}
Regard each proposal as carrying independent uniform marks that randomize the latent kernel and acceptance decision. Conditional on the observable past, its acceptance probability for $e$ is
$\int\widehat\lambda_e(t,h,z,a)\widehat\pi_{t-}(h,\dd z)/c_e$.
The marked-Poisson compensator formula multiplies this probability by the proposal rate $c_e$. Boundedness gives nonexplosion, and \Cref{ass:simulator} identifies the resulting law. The one-stream construction replaces the Bernoulli mark by a categorical mark with the displayed probabilities and the same compensator calculation.
\end{proof}

The kernel is evaluated at the current time even when no market event has occurred. In a finite hidden model with no unobserved transitions between observed jumps, the posterior weights evolve over a quiet interval of length $u$ as
\begin{equation}\label{eq:survival-filter}
 p_s(u)=\frac{p_s(0)e^{-\Lambda_su}}{\sum_vp_v(0)e^{-\Lambda_vu}},
 \qquad\Lambda_s=\sum_e\lambda_e(s).
\end{equation}
After an observed mark, multiply by its conditional rate, apply its hidden-state transition, and normalize. Rejected proposals are internal simulation events and are not added to the observed history. For hidden transitions between events, the exponential in \eqref{eq:survival-filter} is replaced by the appropriate killed transition semigroup.

One exact posterior draw per proposal suffices: it introduces no Monte Carlo rate bias. An approximate posterior or an approximate acceptance rate still incurs the corresponding errors in \Cref{thm:main}. Holding a single draw through an exponential waiting time instead gives a different construction unless a coherent latent model and its filter are maintained. For rates $1$ and $4$ with equal prior weights, the correct latent-mixture survival at time one is
$\tfrac12(e^{-1}+e^{-4})=0.193097540$.
The changing posterior in \eqref{eq:survival-filter} reproduces it; freezing the averaged rate at $2.5$ gives $e^{-2.5}=0.082084999$.

\subsection{The likelihood route and when it is preferable}\label{sec:kl}
Under absolute continuity, the standard point-process likelihood offers a second comparison \citep{bremaud1981}. Assume the same initial law, common event representation, bounded predictable rates, and that the likelihood ratio is valid and integrable. In particular, a strictly positive learned rate on every feasible mark with positive true rate is a convenient sufficient support condition. Then
\begin{equation}\label{eq:path-kl}
 D_{\rm KL}(\mathbb P^a\|\widehat{\mathbb P}^a)
 =\E^a\int_0^T\sum_e\left[
 \bar\lambda_e\log\frac{\bar\lambda_e}{\widehat{\bar\lambda}_e}
 -\bar\lambda_e+\widehat{\bar\lambda}_e\right]\dd t.
\end{equation}
The convention is $0\log(0/q)=0$; a positive true rate against a zero learned rate yields an infinite contribution. Taking the logarithm of the likelihood ratio and replacing each expected counting integral by its compensator proves \eqref{eq:path-kl}. Pinsker's inequality gives the additional bound
\begin{equation}\label{eq:pinsker}
 \TV(\mathbb P^a,\widehat{\mathbb P}^a)
 \le\min\{1,\sqrt{D_{\rm KL}(\mathbb P^a\|\widehat{\mathbb P}^a)/2}\}.
\end{equation}
These are classical likelihood inequalities, not additional novelty claims.

If $\widehat\lambda=(1+\theta)\lambda$ on a set of marks, its KL integrand equals
$\lambda[\theta-\log(1+\theta)]
=\lambda\theta^2/2+O(\lambda|\theta|^3)$.
Thus the KL bound scales as $|\theta|\sqrt{\E\int\lambda}/2$ to first order, whereas the common-event $L^1$ bound scales as $|\theta|\E\int\lambda$. Neither dominates universally. The $L^1$ route handles support mismatch and separates field, posterior, and implementation errors without a denominator containing a small rate. When both apply, use the smaller certificate. A terminal-state TV measurement is not a full-path TV measurement and should not be used to declare either path bound sharp.

\subsection{A mean posterior bound from a coupled latent model}\label{sec:posterior-stability}
The following finite-state observation is useful when the learned filter is derived from a complete simulator, as in \Cref{sec:complete-model}.

\begin{proposition}[Posterior stability in expectation]\label{prop:filter-stability}
Let $P,Q$ be joint laws of observed history $Y$ and a finite hidden state $S$. Let $p(\cdot\mid Y),q(\cdot\mid Y)$ be posterior versions, with $q$ defined on all required $P_Y$ histories. Then
\begin{equation}\label{eq:mean-posterior-tv}
 \E_{P_Y}\TV(p(\cdot\mid Y),q(\cdot\mid Y))\le2\TV(P,Q).
\end{equation}
For a common finite hidden-state model with equal initial law, the same marked transition maps, and total outgoing rate discrepancy bounded by $\epsilon_\lambda$, common-event coupling gives
$\TV(P_t,Q_t)\le1-e^{-\epsilon_\lambda t}$ for the joint history and current hidden state. If the hidden-state metric has diameter $D$, then
\begin{equation}\label{eq:mean-posterior-wone}
 \E_P W_1(\pi_{t-},\widehat\pi_{t-})
 \le2D(1-e^{-\epsilon_\lambda t}).
\end{equation}
\end{proposition}
\begin{proof}
Use a common dominating measure for the observed-history marginals, with densities $a,b$. The triangle inequality gives
$\sum_s|a p_s-a q_s|\le\sum_s|a p_s-b q_s|+|a-b|$.
Integrate and divide by two. The resulting terms are $\TV(P,Q)$ and $\TV(P_Y,Q_Y)$, and marginalization contracts TV. For the dynamic statement couple equal full initial states and use common-event thinning; prior to disagreement the full histories and transitions agree. Finally a maximal coupling on the finite state set gives $W_1\le D\TV$, proving \eqref{eq:mean-posterior-wone}.
\end{proof}

This is an average over true histories, exactly the type of quantity used by \Cref{thm:main}. It is not a uniform pathwise filter bound: rare observed histories can have much larger posterior discrepancies. The $W_1$ form of \Cref{lem:posterior} is preferable here to replacing $W_1$ by $W_2$ and introducing a square-root loss for small finite-state errors.


\subsection{Finite posterior sampling}
Suppose the learned averaged rate is estimated at a fixed history from $N$ independent draws $\widehat Z_i\sim\widehat\pi$. For finite marks,
\begin{equation}\label{eq:mc-rates}
 \E\sum_e\left|
 \frac1N\sum_{i=1}^N\widehat\lambda_e(\widehat Z_i)
 -\E_{\widehat\pi}\widehat\lambda_e(\widehat Z)\right|
 \le \frac1{\sqrt N}\sum_e
 \sqrt{\operatorname{Var}_{\widehat\pi}(\widehat\lambda_e(\widehat Z))}.
\end{equation}
Indeed, each sample mean has variance divided by $N$, and Cauchy--Schwarz bounds its mean absolute error. This is a conditional expectation bound, not a deterministic uniform-in-history bound. Reusing particles or drawing them only at decision times changes their temporal dependence. To apply \Cref{thm:main}, include their randomness in the predictable simulation construction and verify the integrated error under that construction.

\subsection{What resolution independence does and does not imply}
The constant $\sqrt M$ is stable when physical cells are subdivided or quadrature is refined using a fixed reference measure. If instead an unweighted vector norm assigns unit mass to every newly added coordinate, $M$ grows with the representation. The resulting $\sqrt{|\mathcal E|}$ factor is then a change of geometry, not a contradiction.

Neither \Cref{lem:cells} nor \Cref{thm:main} proves that a coarse book path is an autonomous projection of a fine book path. Fine events can change coarse transition probabilities, and hidden reserves can change future displayed behavior. The theorem controls the law generated at the stated physical event resolution. A claim of pathwise consistency across different market resolutions requires an additional compatible transition construction.
